\subsection{Constructing a Bank of Interaction Graphs}
\label{sec:construction}

An experience should explain both what the robot attempted and what its feedback justified doing next. We therefore construct graphs from recorded interactions before building retrieval indices. A graph node represents one interaction phase, such as approaching an object, attempting a grasp, inspecting the result, or placing an object. Objects are arguments of phase predicates, not another type of node. This choice makes the graph describe a decision process whose structure can be reused as a plan.

\paragraph{From recordings to phase intervals.}
The input episode $\tau$ contains a task description, timestamped observations, proposed and executed actions, and controller receipts. We first split source episodes into data partitions, then annotate each episode within its partition. Records retain their sensor, coordinate-frame, action-interface, and source versions. A missing measurement remains unavailable; for example, a commanded gripper closure does not supply measured contact or object retention.

The proposed annotation pipeline has three passes. First, it proposes boundaries at changes in commanded skill or target role, gripper events, execution rejection or interruption, and observable changes relevant to the task. Second, a constrained annotator groups these candidates into intervals and assigns phase labels from a declared vocabulary. Its input includes the surrounding images, action receipts, and candidate event descriptions. Each interval has one immediate objective and a consistent action or observation role. Free-space movement towards one target may span many controller steps; a failed grasp followed by a new inspection creates separate phases because both the objective and the required evidence change. Labels outside the vocabulary are returned for vocabulary revision rather than silently introduced.

Third, validation checks the proposed intervals against their evidence. If a change is only localized between two observations, its boundary is stored as a bracket $[b_i^{-},b_i^{+}]$. A nominal boundary supports indexing, while the stored bracket preserves timing uncertainty. Ambiguous neighboring phases can remain a combined interval until the evidence supports a split. A second annotation pass flags disagreements in phase identity, target association, or outcome. Human adjudication resolves selected disagreements in the training annotations, with the decision and supporting records retained. At deployment, the agent uses the available observation history to estimate its phase. Prefix replay during annotation checks whether a label is already supported at the time it would be used; otherwise that prefix is labeled uncertain.

\paragraph{A common node and evidence schema.}
For $G=(V,E,A)$, the attribute map $A$ assigns each node an identifier, phase label, semantic object-role bindings, interval, preconditions, intended effect, action or skill reference, evidence references, outcome, and provenance. For example, a grasp phase binds the roles \texttt{manipulated\_object} and \texttt{acting\_gripper}. Its intended effect is \texttt{held(object, gripper)}. Whether that effect occurred is a separate field supported by observation references. A place phase uses the same object role and adds \texttt{destination}; its completion requires evidence about the object relative to that destination.

Predicates have typed arguments and a three-valued evidence state, $\operatorname{ver}(p,h_t)\in\{1,0,?\}$, as defined in Sec.~\ref{sec:formulation}.
Each evaluation names its observation or receipt, evaluation procedure, timestamp, and validity conditions. A visual interpretation is marked as estimated, while an execution receipt establishes only the controller facts it reports. Unresolved contradictions produce \(\mathsf{unknown}\). An annotator's confidence is stored separately from predicate truth and is not treated as a calibrated probability. These fields let a model read why a phase succeeded or failed without confusing a requested action with its outcome.

The same schema also represents prospective nodes in $P_t$ and realized nodes in $X_t$. A prospective node specifies a completion predicate and the evidence needed to check it; it has no observed completion value. When executed, it creates a new event node linked to its plan-node identifier and plan version. Historical evidence stays attached to its source episode when the node is retrieved. It cannot certify a predicate in the current scene.

\paragraph{Encoding transitions and recovery.}
Edge records instantiate the readiness rule in Eq.~\ref{eq:readiness}. Each transition or recovery edge records its guard and supporting evidence; each dependency names the required predicate. Mutually exclusive outgoing alternatives share a branch identifier, and at most one can be selected.

Consider putting a bottle in a bin. The first grasp attempt ends with images showing the bottle still on the table after the gripper moves away. Its effect \texttt{held} is false under the declared visual check. A recovery link guarded by that failed effect activates re-observation, which establishes a fresh bottle reference. A second grasp depends on that reference. Once retention is verified, the plan can activate placement and then check that the bottle is inside the bin. The recovery link preserves the relation between the first failed attempt and re-observation; the successful second attempt does not overwrite the failure. If retention is occluded, the graph takes an inspection step instead of assigning success. A retry template specifies a finite attempt budget. Execution unrolls it into distinct event nodes, so repeated attempts preserve their individual actions, feedback, and costs.

\paragraph{Validation and agent-readable serialization.}
The stored representation uses a versioned JSON schema and a predicate dictionary. A validator checks role types, referenced evidence, interval consistency, edge endpoints, branch identifiers, and finite retry budgets. It also checks that a claimed observed effect has a valid evidence reference and that dependencies can be satisfied along an eligible route. Prospective geometry may remain unresolved, but an executable action must name its current grounded target and required capability.

For model input, a deterministic renderer presents each selected node as a short record: role and objective, entry condition, permitted skill, completion evidence, observed outcome if any, and eligible successors. A worked record shows how \texttt{unknown} differs from failure. Both the compact description and the structured record come from the same fields; free-text summaries cannot change the edge guards.

\paragraph{Whole-graph and subgraph indices.}
Each whole graph is indexed by task relation, role types, phase sequence, available capabilities, and entry-state features. A second index stores phase-centered subgraphs. A seed contains a phase and its adjacent transitions; its closure adds the dependency and recovery-trigger records needed to interpret those transitions. Thus, retrieving a regrasp preserves the preceding failure and the observation that made a new attempt reasonable. Each index entry points to immutable graph, node, and media identifiers. Search descriptors represent conditions at the subgraph entry, while the payload retains later outcomes for comparison and adaptation. This separates a queryable situation from the result of the historical response.

\paragraph{Versioned updates.}
Annotation corrections create a new graph version with stable surviving node identifiers and explicit replacements for split or merged phases. Only affected descriptors, subgraph closures, and index entries are rebuilt. During execution, completed or interrupted events extend $X_t$; plan patches create new versions of $P_t$. Neither operation modifies the frozen training bank used in the main evaluation. Cross-episode accumulation instead runs an offline admission step that validates evidence, checks partitions, removes duplicates, and publishes a new bank snapshot. The resulting bank preserves both useful successful procedures and the conditions that motivated recovery.

% CHAPTER_SUMMARY: The proposed annotation pipeline converts source episodes into evidence-linked interaction-phase graphs. Typed conditions, branch semantics, recovery links, stable identities, and validated serialization connect historical experience to prospective plans without treating intended effects as observed outcomes. Whole-graph and dependency-complete subgraph indices preserve the context needed to adapt failures and recoveries; bank updates remain versioned and separate from current-episode execution.
